\begin{claim}{D.73}{Proved}{The full \(q\)-adic cubic quotient and its section kernels}
Retain the actual action, based generators, character and cusp labels
of D.71 at a prime \(q\ne3\) with \(\lambda_q\ne1\).
Write \(\omega=\lambda_q\in\Z_q^\times\), so that
\(\omega^2+\omega+1=0\) and \(\omega-1\) is a unit.
There is an actual quotient
\[
 \begin{gathered}
 Q_\infty=\Z_q\rtimes_\omega C_3,\qquad
 (a,i)(b,j)=(a+\omega^i b,i+j),\\
 \rho_\infty:F\twoheadrightarrow Q_\infty,\qquad
 \rho_\infty(x)=t=(0,1),\quad
 \rho_\infty(y)=r=(1,0),\quad
 \rho_\infty\alpha=\rho_\infty .
 \end{gathered}                                                 \tag{D.73a}
\]
The labelled extension quotient is
\(E/\ker\rho_\infty=Q_\infty\times\langle s\rangle\).
For \(n\geq1\), let \(\rho_n\) be reduction modulo \(q^n\),
and put
\[
 \begin{gathered}
 N_n=\ker\rho_n,\qquad N_\infty=\ker\rho_\infty,\qquad
 C=C_F(s),\\
 C_n=C\cap N_n,\qquad D=C\cap N_\infty .
 \end{gathered}                                                 \tag{D.73b}
\]
The finite-level cover has type
\(((3q^n-1)/2,3)\), and the following are exact:
\[
 \begin{gathered}
 \rho_\infty(C)=\Z_q,\qquad
 H^1(C_n,\F_3)=0\quad(n\geq1),\qquad H^1(D,\F_3)=0,\\
 1\longrightarrow D\longrightarrow C
       \xrightarrow{\ \rho_\infty|_C\ }\Z_q\longrightarrow1 .
 \end{gathered}                                                 \tag{D.73c}
\]
The last sequence splits by a specified procyclic-power construction
in the proof. In particular \(E\) contains the actual subgroup
\(\Z_q\times\langle s\rangle\) supplied by that section.

For both \(n\geq1\) and \(n=\infty\), there are exactly three
\(N_n\)-conjugacy classes of sections in
\(E_n=N_n\rtimes\langle s\rangle\). They have representatives
\[
 T_i=x^i\langle s\rangle x^{-i}\quad(0\leq i<3),
                                                               \tag{D.73d}
\]
and their full normalizers are, respectively,
\(x^i(C_n\times\langle s\rangle)x^{-i}\) and
\(x^i(D\times\langle s\rangle)x^{-i}\).
At each finite level, for the actual rotation projector \(e_n\)
on \(V_n=H^1(N_n,\Z_3)\), the pullback action of
\(\alpha^{-1}\) gives
\[
 (1-e_n)V_n\cong\Z_3[C_3]^{\,q^n-1};
 \qquad e_nV_n\bmod3\cong J_2\oplus J_2 .
                                                               \tag{D.73e}
\]
The first isomorphism is obtained by lifting cyclic vectors,
with the actual action retained.

There is also an exact finite obstruction to two-section generation:
the closed image of every continuous homomorphism
\[
                         C_3\amalg C_3\longrightarrow E
                                                               \tag{D.73f}
\]
is proper. This is proved in the actual quotient
\(Q_1\times C_3\), with both full section classes labelled below.
These statements retain the remaining cubic-coordinate question:
they do not assert existence or exclusion of that coordinate,
or a presentation of the full extension by its section subgroups.
\end{claim}
\begin{proof}
\emph{The quotient before reduction at \(q\).}
The based character change in D.71 keeps
\(c=[x,y]\) exactly and gives
\(\chi(x)=1,\chi(y)=0\). Its index-three cover has
\[
 F_0=\ker\chi,\qquad
 u=x^3,\quad B_i=x^iyx^{-i}\ (0\leq i<3),\qquad
 B_3=uB_0u^{-1}.
                                                               \tag{D.73g}
\]
Use the continuous additive character already constructed
integrally in (D.71l)--(D.71m):
\[
 \ell_q:F_0\longrightarrow\Z_q,\qquad
 \ell_q(u)=0,\quad\ell_q(B_i)=\omega^i,\qquad
                    \ell_q\alpha=\ell_q .
                                                               \tag{D.73h}
\]
In particular its image is all of \(\Z_q\), since it contains \(1\).
Conjugation by \(x\) multiplies it by \(\omega\);
conjugation by \(F_0\) fixes it. Thus its kernel is a
closed normal subgroup of \(F\), preserved in both directions
by the actual automorphism \(\alpha\).

The universal homomorphism from the full free profinite group
which sends \(x\) to \(t\) and \(y\) to \(r\) in (D.73a)
is onto: the closure of the powers of \(r\) is the entire
normal \(\Z_q\), and \(t\) gives the complement.
Its restriction to \(F_0\), followed by the rotation coordinate,
has exactly the four values (D.73h), so its kernel is
\(\ker\ell_q\). The inverse quotient map is
\[
 Q_\infty\longrightarrow F/\ker\ell_q,\qquad
 (a,i)\longmapsto y^{e_q a}x^i\ker\ell_q ,
                                                               \tag{D.73i}
\]
where \(e_q a\in\widehat\Z\) has \(q\)-coordinate \(a\)
and every other primary coordinate zero.
Here \(y^{e_q a}\ker\ell_q\) is the rotation \(a\),
\(x^3\in\ker\ell_q\), and conjugation by \(x\) is
multiplication by \(\omega\). These verify multiplication
and both inverse compositions, including the labels.

The induced automorphism fixes \(r\) and sends \(t\)
to \(r^j t\) for some \(j\in\Z_q\). Since it fixes every
rotation, its actual cube sends \(t\) to \(r^{3j}t\).
The equality \(\alpha^3=1\) and invertibility of \(3\)
in \(\Z_q\) give \(j=0\). This proves
\(\rho_\infty\alpha=\rho_\infty\).
The extension comparison and its inverse are consequently
\[
 \begin{aligned}
 E/N_\infty&\longrightarrow Q_\infty\times C_3,
       &fs^jN_\infty&\longmapsto(\rho_\infty(f),s^j),\\
 Q_\infty\times C_3&\longrightarrow E/N_\infty,
       &(z,s^j)&\longmapsto fs^jN_\infty
                       \quad(\rho_\infty(f)=z).
 \end{aligned}                                                 \tag{D.73j}
\]
The choice of \(f\) changes neither coset, and the maps
are inverse continuous homomorphisms.

\emph{The fixed subgroup and its explicit procyclic section.}
D.48 gives \(H^1(C,\F_3)=0\). Therefore
\(\chi(C)=0\), so \(\rho_\infty(C)\) lies in the rotation
\(\Z_q\). It is closed, since it is a continuous image
of the compact group \(C\). D.71 proves that its
reduction modulo \(q\) is all of \(C_q\).
Every nonzero closed subgroup of the additive group
\(\Z_q\) is \(q^a\Z_q\): choose an element with least
valuation, whose integer powers are dense in that
subgroup. Reduction modulo \(q\) thus forces \(a=0\).
This proves the surjection in (D.73c).

Choose \(h\in C\) with \(\rho_\infty(h)=r\).
For \(a\in\Z_q\) define
\[
                 \iota(a)=h^{e_q a}.
                                                               \tag{D.73k}
\]
Continuous powers in the actual procyclic closure of
\(\langle h\rangle\) show that this is a continuous
homomorphism. Its composite with \(\rho_\infty\) is
\(a\mapsto r^a\), so it is injective and splits the
displayed surjection. In particular the domain is the
full \(\Z_q\), rather than a finite cyclic quotient.
With the conjugation action given by this actual
section there are inverse maps
\[
 \begin{aligned}
 D\rtimes_\iota\Z_q&\longrightarrow C,
                   &(d,a)&\longmapsto d\iota(a),\\
 C&\longrightarrow D\rtimes_\iota\Z_q,
                   &h'&\longmapsto
          \bigl(h'\iota(\rho_\infty(h'))^{-1},
                        \rho_\infty(h')\bigr).
 \end{aligned}                                                 \tag{D.73l}
\]
In these expressions the value of \(\rho_\infty\) on \(C\)
means its additive rotation coordinate.
Every \(\iota(a)\) commutes with \(s\).
Thus \((a,s^j)\mapsto\iota(a)s^j\) gives the asserted
embedded direct product, with inverse given by the
rotation and original quotient-generator labels.

\emph{Every finite power cover, with its integral lattice.}
Put \(m=q^n\), and let \(\omega_n\) denote the reduction
of \(\omega\) modulo \(m\). The quotient \(Q_n\) has
order \(3m\), with exact multiplication obtained by
reducing (D.73a); the actual quotient maps between levels
are reduction of the first coordinate.
The cusp \(c=[x,y]\) has image \(r^{\omega_n-1}\),
which generates the rotation \(C_m\), since
\(\omega-1\) is a \(q\)-adic unit.
It has three cycles of length \(m\) on the \(3m\)
covering sheets. Thus \(N_n\) has three cusps and
Euler characteristic \(-3m\), giving the stated type
\(((3m-1)/2,3)\). A set of actual cusp representatives is
\[
                          d_i=x^i c^m x^{-i}\quad(0\leq i<3).
                                                               \tag{D.73m}
\]

Here the graph-cycle computation in D.71 has the following
literal version for \(m\). Set
\(\Lambda_m=\Z_3[C_m]\), \(\nu_m=1+r+\cdots+r^{m-1}\),
and let \(a,b\in\{1,\ldots,m-1\}\) represent
\(\omega_n,\omega_n^2\). The four Schreier generators
\(u,B_0,B_1,B_2\) have rotation images \(0,1,a,b\).
The graph boundary and its inverse cycle coordinates are
\[
 \begin{aligned}
 \partial(A,B,D',E')&=(r-1)B+(r^a-1)D'+(r^b-1)E',\\
 (A,D',E',k)&\longmapsto
    (A,-P_aD'-P_bE'+k\nu_m,D',E'),\\
 P_j&=1+r+\cdots+r^{j-1}.
 \end{aligned}                                                 \tag{D.73n}
\]
The inverse reads \(A,D',E'\) and the unique \(k\) for
which \(B+P_aD'+P_bE'=k\nu_m\).
Indeed the annihilator of \(r-1\) in \(\Lambda_m\)
is exactly \(\Z_3\nu_m\), by comparing adjacent
coefficients. These are actual inverse
\(\Lambda_m\)-module maps, so
\[
 H_1(N_n,\Z_3)\cong\Z_3\oplus\Lambda_m^3 .
                                                               \tag{D.73o}
\]
Duality gives the same rotation multiplicities for \(V_n\).
The projector \(e_n=m^{-1}\sum_{g\in C_m}g\) is integral,
and restriction and transfer identify \(e_nV_n\) with
the actual rank-four module \(H^1(F_0,\Z_3)\).
They are equivariant for \(\alpha\), since the actual
quotient action is the identity. Its two Jordan-two
blocks are exactly those computed in (D.71s).

For completeness, the paired-lattice argument applies at
every level with the following precise quotient maps.
Fill the three actual cusps (D.73m), then take the
maximal pro-three quotient:
\[
 P_n=\bigl(N_n/
       \overline{\langle\!\langle d_0,d_1,d_2\rangle\!\rangle}
                                                  \bigr)^{(3)} .
                                                               \tag{D.73p}
\]
It is the closed orientable genus-\((3m-1)/2\)
surface pro-three group by the same finite-target
filling universal property used in D.71.
Write the cusp normalizer element as \(v=f_0s\);
then \(\alpha(c)=f_0^{-1}c^\lambda f_0\).
For \(g=x^i\), the quotient \(Q_n/C_m\) selects a unique
\(j\in\{0,1,2\}\); since \(\rho_n(c)\) generates \(C_m\),
there is \(k\in\Z\) and \(h\in N_n\) such that
\[
 \alpha(g)f_0^{-1}=h x^j c^k,\qquad
 \alpha(g c^m g^{-1})=h d_j^\lambda h^{-1}.
                                                               \tag{D.73q}
\]
The inverse automorphism gives the opposite inclusion
of the normal closures. Thus the actual action descends
to (D.73p).
The injection \(H^1(P_n,\Z_3)\to V_n\) has as its image
exactly the characters killing the three \(d_i\).
All cusp evaluations are rotation-invariant, so
\(W_n=(1-e_n)V_n\) lies in this image.
The averaged filled summand has rank two, with the
nonidentity integral action dual to the base torus
matrix \(A_3\); this is the same restriction/transfer
comparison and cusp vanishing test as in D.71,
with \(m\) a unit in \(\Z_3\).
The rotation deck transformations extend positively
over the filled cusps. Hence \(e_n\) is self-adjoint
for the unimodular alternating cup form, and its
orthogonal summand \(W_n\) is integrally unimodular.

Choose a finite unramified extension \(\mathcal O/\Z_3\)
containing all \(m\)-th roots of unity. The idempotents
\[
       e_\psi=\frac1m\sum_{g\in C_m}\psi(g)^{-1}g
                                                               \tag{D.73r}
\]
split \(W_n\otimes\mathcal O\) into rank-three lattices
for all nontrivial \(\psi\).
Since \(m\) is odd, \(\psi\ne\psi^{-1}\); the cup form
pairs precisely those two lattices, perfectly.
The action \(S\), pullback under \(\alpha^{-1}\),
commutes with all these projectors.
Its action on \(H^2(P_n,\Z_3)\) is trivial:
it has order dividing three, whereas \(\Z_3^\times\)
has no nontrivial element of order three.
Thus \(S\) preserves the cup form used here.

The integral fixed-cup lemma proved in D.71 applies
to the actual action on \(P_n\).
Its nonidentity hypothesis, including every inner-twisted
section action, follows from the retained rank-two
averaged action just specified.
Thus pairings of two \(S\)-fixed integral characters
are divisible by three, and the same is true after
flat extension to \(\mathcal O\).
The polynomial \(X^2+X+1\) remains irreducible over
the fraction field of this unramified extension.
On each rank-three character lattice the fixed rank
is therefore one or three. Fixed rank three on
one lattice would give the identity there and on
its perfectly paired inverse lattice, contradicting
that divisibility and the perfect form.
The fixed rank is consequently one on both.
Choose primitive fixed generators \(u,u'\).
For \(w\) with \(\langle w,u'\rangle=1\), write
\[
   (1+S+S^2)w=b u,\qquad
                  b\langle u,u'\rangle=3 .
                                                               \tag{D.73s}
\]
The divisibility and unramified valuation show that
\(\langle u,u'\rangle\) is three times a unit and
\(b\) is a unit. Thus the norm has rank one modulo
three on this rank-three lattice, giving one Jordan
block \(J_3\).
After scalar extension all nontrivial character
summands have this property. Hence over \(\F_3\)
the norm has rank \(m-1\) on a space of dimension
\(3(m-1)\). Every Jordan block has size at most three;
the rank forces every one to have size three.
Lift cyclic vectors for these actual blocks.
The resulting homomorphism
\(\Z_3[C_3]^{m-1}\to W_n\) reduces to an isomorphism
modulo three, hence has unit determinant and an
equivariant inverse. This proves (D.73e) at every
finite level, without a descent or generation
assumption about the full groups.

\emph{The exact section classes and their fixed kernels.}
The other full section generator \(s_1\) has
\(\chi(s_1)=-1\), whereas \(\chi(s)=0\), as in (D.71e).
Consequently none of its \(F\)-conjugates belongs
to \(E_n\), for finite \(n\) or for \(n=\infty\).
Every \(F\)-conjugate of \(s\) does belong, because
the image of \(s\) in (D.73j) is central.
The criterion for two of these sections to be
\(N_n\)-conjugate gives the exact bijections
\[
 \begin{aligned}
 Q_n/\rho_n(C)&\longrightarrow
        \{\text{\(N_n\)-classes of sections in \(E_n\)}\},\\
 \rho_n(f)\rho_n(C)&\longmapsto
                         [f\langle s\rangle f^{-1}] .
 \end{aligned}                                                 \tag{D.73t}
\]
The identical maps apply to \(Q_\infty\).
Indeed a conjugacy between the labelled generators
holds precisely when the corresponding elements
of \(F\) differ by a left \(N_n\)-factor and a
right \(C\)-factor. An \(E_n\)-conjugator may be
multiplied by an element of the source section to
remove its quotient label, so it gives the same
classes as \(N_n\)-conjugacy.
The proved \(\rho_\infty(C)=\Z_q\) makes every coset
space in (D.73t) the three cosets represented by
\(1,t,t^2\).
Any normalizer fixes the labelled generator,
because its conjugation action on \(E/F=C_3\) is
trivial. Removing its quotient power gives exactly
the normalizers stated after (D.73d).

At finite level, the LHS sequence for
\(N_n\to E_n\to C_3\), together with (D.73e), bounds
\[
                    \dim H^j(E_n,\F_3)\leq3\qquad(j\geq2).
                                                               \tag{D.73u}
\]
In fact the two \(J_2\) blocks each contribute one
to positive cyclic-group cohomology, while the
regular blocks contribute zero.
The subgroup \(N_n\) is open with cohomological
dimension one, and every finite subgroup of \(E_n\)
injects into \(C_3\). The full profinite Brown
restriction theorem \cite[Theorem~2 and \S1.1]{A-KochloukovaZalesskii2007}
therefore applies, with no
\(C_3\times C_3\) subgroup. The displayed bound
makes its dense restriction embedding onto.
Each \(C_n\) is closed in \(F\), of
\(3\)-cohomological dimension at most one.
The three actual direct-product normalizers give
\[
        3\bigl(1+\dim H^1(C_n,\F_3)\bigr)
                            =\dim H^j(E_n,\F_3)\leq3
                                  \qquad(j\geq2).
                                                               \tag{D.73v}
\]
This proves the finite-level vanishing.

It also proves the vanishing for \(D\), as follows
directly on continuous characters.
Let \(\xi:D\to\F_3\) be continuous.
Choose an open normal subgroup \(U\) of \(C\) with
\(D\cap U\subseteq\ker\xi\).
Then \(\xi\) descends to \(D/(D\cap U)\cong DU/U\)
and inflates to a character of \(DU\).
The quotient \(C/D=\Z_q\) identifies the open
subgroup \(DU\) with the preimage of some
\(q^n\Z_q\), including possibly \(n=0\).
Thus \(DU=C_n\) when \(n\geq1\), or \(DU=C\)
when \(n=0\). All these character groups vanish,
by (D.73v) and D.48. Hence \(\xi=0\).
The normalizer and class statements at \(n=\infty\)
were already proved by the actual coset maps;
no finite-type surface assertion for \(N_\infty\)
is made.

\emph{The finite image of any two torsion generators.}
Reduce (D.73j) modulo \(q\), writing
\(Q_1=\langle r,t\rangle\).
Every nonidentity finite-order element of \(E\)
has order three, since \(F\) is torsion-free and
its image in \(E/F\) is injective.
Replace it by its inverse if needed to give the
specified quotient-generator label \(s\).
The two full section classes then have images
\[
                  s,\qquad r^a t^{-1}s\quad(a\in\F_q).
                                                               \tag{D.73w}
\]
For the first class this follows from centrality
of \(s\) in the quotient; for the other it follows
from \(\chi(s_1)=-1\), and conjugation by rotations
gives exactly the displayed possibilities because
\(\omega_1-1\) is a unit.

Two elements from the first class generate at most
\(\langle s\rangle\).
One element from each class generates the order-nine
abelian subgroup
\(\langle s,r^a t^{-1}s\rangle\), which omits the
order-\(q\) rotation.
Two from the second class lie in the kernel of
the explicit homomorphism
\[
 Q_1\times C_3\longrightarrow C_3,\qquad
                       (r^a t^i,s^j)\longmapsto i+j .
                                                               \tag{D.73x}
\]
This is again a proper subgroup. The conclusions
also hold if one image is the identity.
Thus every subgroup generated by two finite-order
elements has proper image in this particular
finite quotient. A homomorphism from the full
free profinite product in (D.73f) has precisely
such generator images, proving its proper-image
assertion. No converse generation statement is used.

The exact producer
\path{code/part_c/torus_cubic_qadic_d73.py}
checks the compatible prime-power roots, quotient
coordinate maps, three cusp cycles, graph-module
inverse formulas, procyclic primary splitting
coordinates, and all two-section image types in
the specified finite cases.
It passes 155,189 exact checks over fourteen
prime-power cases and 3,252 section-image pairs.
The infinite quotient, fixed-kernel vanishing
and integral lattice conclusions are the actual
maps and cohomology arguments above.
\end{proof}
